\subsection{\texorpdfstring{AFFINE LINEAR VARIETIES IN \(R^{n}\)}{AFFINE LINEAR VARIETIES IN R TO THE N}}

Given a \(p\) -dimensional affine linear variety \(\mathbf{V}\) in \(\mathbf{R}^n\) , there exists an affine mapping \(f\) of \(\mathbf{R}^n\) onto itself which transforms \(\mathbf{V}\) into a \(p\) -dimensional coordinate variety, that is to say a vector subspace \(\mathbf{V}'\) generated by \(p\) of the vectors of the canonical basis \((\boldsymbol{e}_i)\) of \(\mathbf{R}^n\) . There exists a mapping of \(\mathbf{V}'\) onto \(\mathbf{R}^p\) which is an isomorphism of the vector space structure and the topology of \(\mathbf{V}'\) onto the corresponding structures of \(\mathbf{R}^p\) (it is moreover often convenient to identify \(\mathbf{R}^p\) with such a coordinate variety \(\mathbf{V}'\) , e.g., with the vector subspace generated by \(\boldsymbol{e}_1, \boldsymbol{e}_2, \ldots, \boldsymbol{e}_p\) ). In addition, \(\mathbf{V}'\) is a closed subset of \(\mathbf{R}^n\) (Chapter I, § 4, no. 3, Corollary to Proposition 7); hence:

PROPOSITION 2. Every p-dimensional linear affine variety in \(\mathbf{R}^n\) is a closed subset of \(\mathbf{R}^n\) , homeomorphic to \(\mathbf{R}^p\) .

It is this result which is the origin of the names line and plane given to affine linear varieties of one and two dimensions in a vector space over an arbitrary division ring. We recall also that, for \(n \geqslant 1\) , the affine linear varieties of n - 1 dimensions of \(R^{n}\) are called hyperplanes (loc. cit.).

The n one-dimensional coordinate varieties, that is to say the n lines passing through o and the n points \(e_{i}\) respectively, are called the coordinate axes. For n = 2 the axis through \(e_{1}\) is sometimes called the axis of abscissas and the axis through \(e_{2}\) is called the axis of ordinates; the first coordinate of a point \(x \in R^{2}\) is called its abscissa, the second coordinate its ordinate.

Every line D passing through a point \(\pmb{a}\) has a parametric representation \(t \to \pmb{a} + tb\) , where \(t\) runs through \(\mathbf{R}\) and \(b \neq 0\) ; this mapping is a homeomorphism of \(\mathbf{R}\) onto \(D\) . The vector \(\pmb{b}\) is called a direction vector of D, and its components \(b_{i}\) ( \(1 \leqslant i \leqslant n\) ) are called direction ratios of D. If \(b'\) is another direction vector of D, there exists \(h \neq o\) in R such that \(b' = hb\) .

The set of points \(a + tb\) , where t runs through the set of real numbers \(\geqslant 0\) , is called the closed ray (or simply ray, or half-line) with origin a and direction vector b (or with direction ratios \(b_{i}\) ). It is a closed subset of \(R^{n}\) , homeomorphic to the interval \([0, +\infty[\) of R, and therefore connected. The line D is the union of the two rays with origin a and direction vectors b and -b respectively; these rays are said to be opposite.

By abuse of language, the set of points \(\pmb{a} + t\pmb{b}\) , where \(t\) runs through the set of real numbers \(> 0\) , is called the open ray with origin \(\pmb{a}\) and direction vector \(\pmb{b}\) ; it is homeomorphic to the interval \(]0, +\infty[\) (and therefore homeomorphic to \(\mathbf{R}\) ), but is not open in \(\mathbf{R}^n\) if \(n > 1\) , although it is open in the line which contains it.

A line passing through two distinct points x and y also has a parametric representation \((u, v) \to ux + vy\) , where \((u, v)\) runs through the set of pairs of real numbers such that \(u + v = 1\) . Given any two points x, y (distinct or not) the set of points \(ux + vy\) , where \((u, v)\) runs through the set of pairs of real numbers such that \(u \geqslant 0\) , \(v \geqslant 0\) and \(u + v = 1\) , is called the closed segment (or simply segment) with end-points x, y. A closed segment is compact and connected, for if its end-points are distinct it is homeomorphic to the interval {[}0, 1{]} of R.

If \(x \neq y\) , the set of points \(ux + vy\) such that \(u > 0, v > 0\) and \(u + v = 1\) is called (by abuse of language) the open segment with end-points \(x, y\) ; it is homeomorphic to the open interval \(]0, 1[\) (and hence also homeomorphic to \(R\) ). Finally the union of \(\{y\}\) and the open segment with endpoints \(x, y\) is sometimes called the segment open at \(x\) and closed at \(y\) ; it is homeomorphic to the interval \([0, 1[\) . All the segments with \(x\) and \(y\) as end-points are connected, and the closure of each of them is the closed segment with the same end-points.

PROPOSITION 3. Let \(f(x) = f(x_1, x_2, \ldots, x_n)\) be a polynomial with real coefficients, not identically zero, defined on \(\mathbf{R}^n\) . Then the complement of the set \(\overline{f}(0)\) is dense in \(\mathbf{R}^n\) .

Let \(\pmb{x}\) be any point of \(\mathbb{R}^n\) and let \(y \in \mathbb{R}^n\) be such that \(f(y) \neq 0\) ; \(\varphi(t) = f(x + t(y - x))\) is a polynomial in the real variable \(t\) , not identically zero; hence there exist arbitrarily small values of \(t\) such that \(\varphi(t) \neq 0\) . This shows that \(\pmb{x}\) lies in the closure of the complement of \(\overline{f}(0)\) .

COROLLARY. The complement of an affine linear variety of dimension \(p < n\) is dense in \(\mathbf{R}^n\) .

Since every affine linear variety of dimension p \textless{} n is contained in a hyperplane, it is enough to prove the corollary for a hyperplane; but a hyperplane is defined by an equation \(g(x) = 0\) , where g is a linear polynomial not identically zero.

PROPOSITION 4. In \(\mathbf{R}^n\) ( \(n \geqslant 1\) ) the complement of every hyperplane has two connected components.

Let \(g(x) = 0\) be an equation of a hyperplane \(H\) in \(\mathbb{R}^n\) , \(g\) being a linear polynomial. The set \(\complement H\) is the union of the set \(E_1\) of all points \(x\) such that \(g(x) > 0\) and the set \(E_2\) of all points \(x\) such that \(g(x) < 0\) . \(E_1\) and \(E_2\) are connected, for if \(g(x) > 0\) and \(g(y) > 0\) we have \(g(ux + vy) = ug(x) + vg(y) > 0\) whenever \(u \geqslant 0\) , \(v \geqslant 0\) and \(u + v = 1\) ; in other words, the segment with end-points \(x\) and \(y\) is contained in \(E_1\) . Similarly for \(E_2\) . On the other hand, \(\complement H\) is not connected, because its image in \(R\) under \(g\) is the union of the intervals \(]0, +\infty[ \text{and}] - \infty, o[\) .

The components \(E_{1}\) , \(E_{2}\) of the complement of a hyperplane H are called the open half-spaces determined by H.

The closures of \(E_{1}\) and \(E_{2}\) , which are respectively \(E_{1} \cup H\) and \(E_{2} \cup H\) , are called the closed half-spaces determined by H.

Observe that an affine mapping of \(R^{n}\) onto itself which transforms H into a ``coordinate'' hyperplane, e.g., the hyperplane whose equation is \(x_{n}=0\) , also transforms the open half-spaces determined by H into the open half-spaces defined respectively by the relations \(x_{n}>0\) and \(x_{n}<0\) ; the latter are open boxes and therefore homeomorphic to \(R^{n}\) .

